% The GNS quotient, exercised on a genuinely degenerate reflection pairing.
% Lane: O (operator bridge), modules YangMills/OS/**.
% REGIME: tricotomy; NO scores, NO desk language, NO commit chronology beyond
% the reproducibility freeze.  Freeze: Lean tree of source checkpoint
% efc42c3179827b9985838c86782dcc05b5f9a529, full core build 8423 jobs,
% 2451 oracle commands (see Section 7).
\documentclass[11pt]{article}
\usepackage[margin=1.1in]{geometry}
\usepackage{amsmath,amssymb,amsthm}
\usepackage{booktabs,longtable,array}
\usepackage[colorlinks=true,linkcolor=blue,urlcolor=blue,citecolor=blue]{hyperref}

\newtheorem{theorem}{Theorem}[section]
\newtheorem{lemma}[theorem]{Lemma}
\newtheorem{proposition}[theorem]{Proposition}
\newtheorem{corollary}[theorem]{Corollary}
\newtheorem{remark}[theorem]{Remark}
\theoremstyle{definition}
\newtheorem{definition}[theorem]{Definition}

\emergencystretch=3em
\tolerance=2000

\newcommand{\lean}[1]{\texttt{#1}}
\newcommand{\anchor}{efc42c3179827b9985838c86782dcc05b5f9a529}
\newcommand{\osline}[3]{%
  \href{https://github.com/lluiseriksson/THE-ERIKSSON-PROGRAMME/blob/\anchor/YangMills/OS/#1.lean\#L#2}{\nolinkurl{#3}}}
\newcommand{\repofile}[2]{%
  \href{https://github.com/lluiseriksson/THE-ERIKSSON-PROGRAMME/blob/\anchor/#1}{\nolinkurl{#2}}}
\newcommand{\R}{\mathbb{R}}
\newcommand{\C}{\mathbb{C}}
\newcommand{\Z}{\mathbb{Z}}
\newcommand{\SU}{\mathrm{SU}}
\newcommand{\Om}{\Omega}
\newcommand{\inn}[2]{\langle #1,\, #2\rangle}
\newcommand{\Ex}{\mathbb{E}}

\title{The Quotient That Is Not the Identity:\\
A Machine-Checked Degenerate Reflection Pairing\\
and Its Gelfand--Naimark--Segal Quotient\\
\large Four observables, two states, and the collapse in between}
\author{Lluis Eriksson}
\date{}

\begin{document}
\maketitle

\begin{abstract}
The Osterwalder--Seiler reconstruction passes from a reflection-positive measure
to a Hilbert space by quotienting out the null space of the reflected pairing.
In a companion development that step was present but did nothing: the pairing
there was definite, so the quotient was the identity, and that paper says so in
its own abstract.  This paper supplies the missing case, in Lean~4 with
\texttt{mathlib}.

For the $\Z_2$ lattice gauge chain we take half-space observables of two time
slices --- a four-dimensional space --- and form the reflected pairing directly
from the Boltzmann weights.  The physical space of the chain is two-dimensional.
Integrating out the future collapses four observables onto two states, and that
collapse is the null space.  We prove: the pairing factors through an
independently defined reconstruction map $\Phi$; its self-pairing rearranges
into a manifest sum of two non-negative terms, from which positivity and the
null space follow together; the null space is \emph{exactly} $\ker\Phi$, not
merely non-empty; an explicit non-zero observable lies in it; and the quotient is
isomorphic to the physical space \emph{by the map $\Phi$ itself}, not by a
dimension count.

The degeneracy is the mechanism the reconstruction exists to handle, and it is
the one that reappears at every larger system.  What is not claimed: this is two
time slices and not $m$; still $\Z_2$, one variable per slice, fixed finite
size, and not volume-uniform; $\Z_N$ for $N>2$ is untouched; and the completion
step of the reconstruction is trivial here because every space in sight is
finite-dimensional, which we state rather than present as work done.  Nothing in
this paper is a claim about $\SU(N)$, the continuum limit, or the Yang--Mills
mass gap.
\end{abstract}

\section{Introduction}

\subsection{The step that usually does nothing in small examples}

Reflection positivity gives a positive \emph{semi}-definite pairing on
half-space observables.  The Hilbert space of the reconstruction is obtained by
quotienting out its null space and completing \cite{OsterwalderSeiler,
GlimmJaffe}.  In examples small enough to be fully verified, the pairing is
often already definite, and then the quotient is the identity map --- the step
is performed vacuously.  A formalisation in that situation has not exercised the
construction; it has avoided it.

That was the position of a companion development \cite{ErikssonOSChain}, which
verified an end-to-end reconstruction for the $\Z_2$ chain and stated in its
abstract that its own quotient did no work.  This paper removes that limitation
by exhibiting, for the same measure, a pairing whose null space is provably
non-zero, computing that null space exactly, and identifying the quotient.

\subsection{Where the degeneracy comes from}

Not from a contrivance.  Take half-space observables depending on \emph{two}
time slices, $A : \Z_2 \times \Z_2 \to \C$.  That space is four-dimensional.
The physical space of the chain is two-dimensional.  A four-dimensional space
cannot inject into a two-dimensional one, so the pairing \emph{must} be
degenerate; the content is to compute the degeneracy rather than to observe that
it exists.  Physically, the collapse is the statement that the future has been
integrated out: two observables that agree after that integration are the same
state.

This is the mechanism, and it does not go away at larger systems.  Any lattice
with spatial extent has half-space observables in vastly greater number than the
states of its transfer operator.

\subsection{What is proved}

Write $k_\beta(i,j) = e^{\beta s(i,j)}$ for the Boltzmann bond weight, $s=+1$ on
equal spins and $-1$ otherwise.  The reflected pairing on two-slice observables
is the explicit finite sum
\[
  \inn{A}{B} \;=\; \sum_{a,b,c,d}
     \overline{A(b,a)}\; B(c,d)\; k_\beta(a,b)\,k_\beta(b,c)\,k_\beta(c,d),
\]
the four spins being $\sigma_{-2},\sigma_{-1},\sigma_0,\sigma_1$ with the
reflection carrying the half-space $\{0,1\}$ to $\{-1,-2\}$.  Separately we
define
\[
  (\Phi A)(c) \;=\; \sum_d k_\beta(c,d)\, A(c,d).
\]
Then:
\begin{itemize}
\item $\inn{A}{B} = \sum_{b,c}\overline{(\Phi A)(b)}\,k_\beta(b,c)\,(\Phi B)(c)$
      --- the pairing factors through $\Phi$;
\item $\inn{A}{A} = (e^\beta-e^{-\beta})\bigl(|v_0|^2+|v_1|^2\bigr)
      + e^{-\beta}|v_0+v_1|^2$ with $v=\Phi A$, a manifest sum of two
      non-negative terms;
\item hence $\inn{A}{A}\ge 0$ for $\beta\ge 0$, and for $\beta>0$,
      $\inn{A}{A}=0 \iff \Phi A = 0$;
\item an explicit $A\neq 0$ with $\Phi A = 0$;
\item and $\bigl(\C^{2\times 2}\bigr)/\ker\Phi \cong \C^2$, the isomorphism
      induced by $\Phi$.
\end{itemize}

\subsection{What is deliberately not claimed}

\begin{enumerate}
\item \textbf{Two slices, not $m$.}  The construction is carried out at depth
two.  Depth $m$ is not treated.
\item \textbf{Still the smallest gauge group and no spatial extent.}  One
$\Z_2$ variable per slice; the physical space stays two-dimensional.
\item \textbf{Fixed finite size, not volume-uniform.}  No statement here is
uniform in a volume, and at this size the question does not arise.
\item \textbf{The completion is trivial.}  Every space here is
finite-dimensional, so the completion half of the reconstruction is vacuous.  We
state that; we do not present it as work done.
\item \textbf{Nothing about $\SU(N)$, the continuum limit, or the Clay problem}
\cite{Clay}.  No reduction from any statement in this paper to the Yang--Mills
mass gap is claimed, and none is known to the author.
\end{enumerate}

\section{The pairing, and the map}

\begin{definition}[Reflected pairing]
For $A,B : \Z_2\times\Z_2 \to \C$ set
$\inn{A}{B}$ to be the four-fold sum displayed above.
\end{definition}

In Lean this is \osline{Z2Quotient}{87}{reflPairing}, built from
\osline{Z2Identification}{76}{z2Bond} and from nothing else.

\begin{remark}[Non-circularity]\label{rem:noncirc}
The definition of the pairing mentions no reconstruction map, no transfer
operator, and neither coefficient of the normalised transfer matrix.  Printing
it exhibits its whole dependency set.  This is the condition that makes
Theorem~\ref{thm:factor} a theorem about two independently defined objects
rather than a restatement of a definition; the same discipline governs the
companion development \cite{ErikssonOSChain}.
\end{remark}

\begin{definition}[Reconstruction map]
$(\Phi A)(c) = \sum_d k_\beta(c,d) A(c,d)$, defined without reference to the
pairing: \osline{Z2Quotient}{96}{slicePhi}, and as a linear map
\osline{Z2Quotient}{106}{slicePhiLin}.
\end{definition}

\begin{theorem}[The pairing factors through $\Phi$]\label{thm:factor}
$\displaystyle \inn{A}{B} \;=\; \sum_{b,c}
   \overline{(\Phi A)(b)}\; k_\beta(b,c)\; (\Phi B)(c)$.
\end{theorem}

\begin{proof}[Proof, as formalised]
Group the four-fold sum as $\sum_{b,c}\bigl[\sum_a \overline{A(b,a)}k_\beta(a,b)\bigr]
k_\beta(b,c)\bigl[\sum_d k_\beta(c,d)B(c,d)\bigr]$.  The inner left bracket is
$\overline{(\Phi A)(b)}$ because $k_\beta$ is real and symmetric; the inner right
bracket is $(\Phi B)(c)$.
\end{proof}

Artifact: \osline{Z2Quotient}{128}{reflPairing\_eq}.

\section{One identity, two consequences}

\begin{theorem}[Sum of two non-negative terms]\label{thm:sos}
With $v = \Phi A$,
\[
  \inn{A}{A} \;=\;
  \bigl(e^{\beta}-e^{-\beta}\bigr)\bigl(|v_0|^2+|v_1|^2\bigr)
  \;+\; e^{-\beta}\,|v_0+v_1|^2 .
\]
\end{theorem}

\begin{proof}
Expand Theorem~\ref{thm:factor} over the two-element index set, with
$k_\beta(i,i)=e^\beta$ and $k_\beta(i,j)=e^{-\beta}$ for $i\neq j$, and complete
the square in $v_0+v_1$.
\end{proof}

Artifact: \osline{Z2Quotient}{143}{reflPairing\_self}.  Both facts we need come
from this one line.

\begin{corollary}[Reflection positivity]\label{cor:pos}
$\inn{A}{A}\ \ge\ 0$ for every $\beta\ge 0$.
\end{corollary}

\osline{Z2Quotient}{159}{reflPairing\_self\_nonneg}.

\begin{corollary}[The null space, exactly]\label{cor:null}
For $\beta>0$: $\;\inn{A}{A}=0 \iff \Phi A = 0$.
\end{corollary}

\osline{Z2Quotient}{171}{reflPairing\_self\_eq\_zero\_iff}.

\begin{remark}[Why $\beta>0$, said rather than avoided]
At $\beta=0$ the bond weight is constant and the kernel $k_\beta$ itself
degenerates, so the coefficient $e^\beta-e^{-\beta}$ of
Theorem~\ref{thm:sos} vanishes and the equivalence of
Corollary~\ref{cor:null} fails.  Positivity survives at $\beta=0$;
definiteness does not.  The hypotheses in the formalisation record exactly this
split rather than excluding the point silently.
\end{remark}

\section{The null space is not zero}

This is the point of the paper: without it the construction would have
reproduced the trivial case.

\begin{theorem}[Explicit null vector]\label{thm:null}
Let $N(c,0) = k_\beta(c,1)$ and $N(c,1) = -k_\beta(c,0)$.  Then $\Phi N = 0$ and
$N \neq 0$; consequently $\inn{N}{N} = 0$ for $\beta > 0$ and the null space is
non-trivial.
\end{theorem}

\begin{proof}
$(\Phi N)(c) = k_\beta(c,0)k_\beta(c,1) - k_\beta(c,1)k_\beta(c,0) = 0$, exactly.
$N\neq 0$ because $N(0,0) = k_\beta(0,1) = e^{-\beta} > 0$: the Boltzmann weights
are strictly positive, which is the only property of them used here.
\end{proof}

Artifacts: \osline{Z2Quotient}{210}{nullObs},
\osline{Z2Quotient}{220}{slicePhi\_nullObs},
\osline{Z2Quotient}{229}{nullObs\_ne\_zero},
\osline{Z2Quotient}{235}{reflPairing\_nullObs}, and as a submodule statement
\osline{Z2Quotient}{273}{nullSubmodule\_ne\_bot}.

\section{The quotient is the physical space}

\begin{theorem}[Quotient]\label{thm:quot}
$\Phi$ is surjective, and the induced map
\[
  \bigl(\Z_2\times\Z_2 \to \C\bigr)\;/\;\ker\Phi \;\;\xrightarrow{\ \sim\ }\;\; \C^2
\]
is a linear isomorphism.  For $\beta>0$ the kernel is exactly the null space of
the pairing, so this is the Gelfand--Naimark--Segal quotient of the reflected
pairing.
\end{theorem}

Artifacts: \osline{Z2Quotient}{252}{slicePhi\_surjective},
\osline{Z2Quotient}{265}{nullSubmodule},
\osline{Z2Quotient}{268}{mem\_nullSubmodule\_iff},
\osline{Z2Quotient}{283}{quotEquivPhysical}.

\begin{remark}[Induced by the map, not by a count]
Four and two are the dimensions, and a reader could observe that a
two-dimensional quotient of a four-dimensional space by a two-dimensional
subspace is $\C^2$ for trivial reasons.  That is not what is proved.  The
isomorphism here is the one \emph{induced by $\Phi$}, so the quotient is the
physical space \emph{because of what the reconstruction map does to
observables}.  A dimension count would leave the identification unanchored to
the physics.
\end{remark}

\begin{remark}[What this adds to the companion paper]
\cite{ErikssonOSChain} verified the chain
measure $\to$ reflection positivity $\to$ Hilbert space $\to$ transfer operator
$\to$ identification $\to$ decay, and proved that \emph{its} quotient step was
the identity.  That statement was true of that system and is unaffected.  The
present paper does not amend it; it supplies the case that paper did not cover.
The two together exercise both branches: a pairing where the quotient is the
identity, with that proved, and a pairing from the same measure where it is not.
\end{remark}

\section{Formalisation notes}

Two decisions are worth recording.

\paragraph{Disjoint definition sets, again.}  The measure side and the map side
share nothing; they meet only in Theorem~\ref{thm:factor}.  This is what makes
Remark~\ref{rem:noncirc} checkable by printing a definition, and it is the same
architecture as the companion development.

\paragraph{Pattern matching instead of conditionals.}  The witnesses of
Theorems~\ref{thm:null} and~\ref{thm:quot} were first written with a conditional
on the two-element index.  A conditional on a decidable equality normalises to a
form that the obvious simplification lemma does not match, and three routine
proofs failed with goals displaying an unreduced conditional.  Defining the
witnesses by pattern matching instead makes their two values reduce
definitionally, and the proofs become rewrites backed by reflexivity.  The
general lesson --- remove the awkward term from the statement rather than defeat
it inside the proof --- is the same one that governs the algebraic steps of the
companion development.

\section{Reproducibility}\label{sec:repro}

All artifacts are at source checkpoint \texttt{\anchor} on branch \texttt{main}
of \url{https://github.com/lluiseriksson/THE-ERIKSSON-PROGRAMME}.  Every link in
this paper was resolved against that commit, and checked to point at an actual
declaration line, before compilation.

\begin{center}
\begin{tabular}{ll}
\toprule
Quantity & Value \\
\midrule
Full core build & 8423 jobs, success \\
Oracle commands & 2451 \\
\quad with axiom dependencies & 2429 \\
\quad axiom-free & 22 \\
Distinct axioms across all outputs & \texttt{propext}, \texttt{Quot.sound}, \texttt{Classical.choice} \\
Occurrences of \texttt{sorryAx} & 0 \\
Project axioms & 0 \\
Errors & 0 \\
\bottomrule
\end{tabular}
\end{center}

The criteria this work was required to meet --- including the non-circularity
condition of Remark~\ref{rem:noncirc} and the requirement that the null space be
proved non-zero --- were fixed in
\repofile{docs/O-BRIDGE-CHARTER.md}{docs/O-BRIDGE-CHARTER.md}, Amendment~9, in a
commit preceding the existence of the module they govern.  Their verdicts are in
\repofile{docs/VERIFICATION-LEDGER.md}{docs/VERIFICATION-LEDGER.md},
Addendum~516.  The oracle script is
\repofile{oracle\_check.lean}{oracle\_check.lean}.

\section{Continuation}

The next step is spatial extent: a transfer operator on a space that grows with
the volume, where the half-space algebra is larger still and the quotient of
this paper becomes indispensable rather than merely non-trivial.  That is where
the present style of construction is expected to stop being easy, and it is the
first step after which volume-uniformity --- the only property in this
neighbourhood that could eventually bear on a mass gap --- becomes a meaningful
question.  It is not reached here.

\begin{thebibliography}{9}

\bibitem{OsterwalderSeiler}
K.~Osterwalder and E.~Seiler,
\emph{Gauge field theories on a lattice},
Annals of Physics \textbf{110} (1978), 440--471.

\bibitem{GlimmJaffe}
J.~Glimm and A.~Jaffe,
\emph{Quantum Physics: A Functional Integral Point of View},
2nd ed., Springer, 1987.

\bibitem{ErikssonOSChain}
L.~Eriksson,
\emph{From the Gibbs Weight to the Spectral Gap: A Complete Machine-Checked
Osterwalder--Seiler Chain for the $\Z_2$ Lattice Gauge Chain}, 2026.

\bibitem{Lean4}
L.~de~Moura and S.~Ullrich,
\emph{The Lean 4 theorem prover and programming language},
CADE-28, Lecture Notes in Computer Science \textbf{12699} (2021), 625--635.

\bibitem{Mathlib}
The mathlib Community,
\emph{The Lean mathematical library},
CPP 2020, 367--381.

\bibitem{Clay}
A.~Jaffe and E.~Witten,
\emph{Quantum Yang--Mills theory},
Clay Mathematics Institute Millennium Prize Problem description, 2000.

\end{thebibliography}

\end{document}
